% !TeX program = pdflatex
% ─────────────────────────────────────────────────────────────────────────────
%  papers/dualsort.tex — O DUAL SORT, em paper próprio.
%
%  Compila sozinho:  pdflatex dualsort.tex
%
%  Nada aqui vem de fora. As peças são a estaca, a Lei 2, a bidualidade e o
%  trial, e todas estão demonstradas na Teoria e medidas em tests/dualsort.c.
% ─────────────────────────────────────────────────────────────────────────────
\documentclass[11pt,a4paper]{article}
\usepackage[utf8]{inputenc}\usepackage[T1]{fontenc}\usepackage[brazil]{babel}
\usepackage{amsmath,amssymb,amsthm}
\usepackage[margin=2.4cm]{geometry}
\usepackage{booktabs}\usepackage{xcolor}
\usepackage[hidelinks]{hyperref}

\definecolor{cinza}{gray}{0.35}
\newcommand{\code}[1]{\texttt{\small #1}}
\newcommand{\medido}[1]{\par\medskip\noindent{\small\color{cinza}\textbf{Medido.} #1}\par\medskip}
\newtheorem{teorema}{Teorema}
\newtheorem{definicao}[teorema]{Definição}
\renewcommand{\proofname}{Demonstração}

\title{\textbf{O Dual Sort}\\[2pt]
  {\large ordenar é descer pela estaca}}
\author{Aarão Melo Lopes}
\date{6 de agosto de 2026}

\begin{document}
\maketitle

\begin{abstract}
\noindent Apresenta-se um algoritmo de ordenação construído inteiramente a partir da teoria da
dualidade: a \textbf{estaca} $x\mapsto x^{\dagger}$, a \textbf{Lei 2} $-f=f^{-1}$, a
\textbf{bidualidade} $J^{4}=\mathrm{id}$ e o \textbf{trial}. O algoritmo não compara elementos
entre si --- o lado de cada elemento sai do \emph{sinal da estaca}, isto é, de que lado da
fronteira ele está --- e a sua terminação não é uma propriedade a verificar: é o fecho da órbita.
Mostra-se ainda que retirar qualquer das peças quebra o algoritmo, e que retirar o ponto fixo o
faz \emph{não terminar}.
\end{abstract}

\section{As peças}

Todas as peças estão demonstradas na \emph{Teoria} e nenhuma é importada.

\begin{definicao}[a estaca]
Num conjunto de centro $c$, a \textbf{estaca} é
\[
  x^{\dagger} \;=\; 2c - x .
\]
Ela é \textbf{involutiva}: $x^{\dagger\dagger}=x$. A sua \textbf{fronteira} é o conjunto dos seus
pontos fixos, $x^{\dagger}=x$, isto é, $x=c$ --- e ele é um só.
\end{definicao}

\noindent A estaca troca de lado. O que ela fixa é o centro, e é por isso que ele não pertence a
nenhum dos lados: \emph{é a fronteira entre eles}.

\begin{definicao}[a Lei 2 e o $J$]
A Lei 2 escreve-se, onde o dual é a inversa,
\[
  -f \;=\; f^{-1},
\]
o que força $f^{2}=-1$. Em duas coordenadas isso é $J(a,b)=(-b,a)$: uma rotação de um quarto.
\end{definicao}

\begin{teorema}[a bidualidade fecha]
$J^{2}=-1$ e $J^{4}=\mathrm{id}$. A órbita de $J$ visita \textbf{quatro} quadrantes distintos e
volta ao primeiro.
\end{teorema}

\begin{proof}
Aplicando $J$ duas vezes a $(a,b)$ obtém-se $(-a,-b)$, que é $-1$ vezes o par; aplicando quatro,
$(a,b)$. Como $J$ leva cada quadrante ao seguinte e $J^{4}$ é a identidade, os quatro são
distintos e o quinto passo é o primeiro. \end{proof}

\noindent \textbf{É este fecho que garante a terminação}, e não uma condição imposta por fora.

\section{O algoritmo}

Seja $A$ uma lista finita. Se $|A|\leq 1$, ela já está ordenada. Caso contrário, tome-se
$c=\min A+\lfloor(\max A-\min A)/2\rfloor$ e classifique-se cada elemento pelo \textbf{sinal da
estaca}:
\[
  x^{\dagger}>x \iff x<c, \qquad
  x^{\dagger}<x \iff x>c, \qquad
  x^{\dagger}=x \iff x=c .
\]

\noindent \textbf{São três, e não dois} --- é o \emph{trial}. A esquerda e a direita descem outra
vez; a fronteira fica entre elas, e não desce, porque não há para onde: ela é o ponto fixo.

\begin{center}\small
\begin{tabular}{@{}ll@{}}
\toprule
o que a peça faz & no algoritmo \\
\midrule
a \textbf{estaca} é involutiva & garante que os dois lados são imagem um do outro \\
a \textbf{fronteira} é o ponto fixo & é ela que \emph{parte}, e não pertence a nenhum lado \\
o \textbf{trial} & três destinos: esquerda, fronteira, direita \\
a \textbf{bidualidade} & o fecho da órbita, e portanto a terminação \\
\bottomrule
\end{tabular}
\end{center}

\subsection*{Não se compara nada}

O destino de um elemento sai do \textbf{sinal de $x^{\dagger}-x$} --- uma leitura sobre o próprio
elemento --- e nunca de uma pergunta sobre outro elemento. \emph{Em por-unidade não há grandeza a
comparar: há posição a ler}, e ela já está no elemento.

\section{O que se mede}

\medido{a estaca é involução em $101$ pontos e o seu ponto fixo é \textbf{um só}; ela parte em
$50$ e $50$, e o dual de cada elemento cai sempre no outro lado; o $J$ visita os quatro quadrantes,
todos distintos, e volta ao quarto; $J^{2}=-1$ e $J^{4}=\mathrm{id}$ em $169$ pares;
\textbf{$200$ listas} saem crescentes e completas descendo pela estaca (\code{tests/dualsort.c}).}

\subsection*{E o que acontece quando se tira uma peça}

Três mutações, e nenhuma delas degrada o resultado: \textbf{quebram-no}.

\begin{center}\small
\begin{tabular}{@{}ll@{}}
\toprule
tirar & o que acontece \\
\midrule
o ponto fixo vai para um dos lados & a partição deixa de reduzir --- \textbf{não termina} \\
o sinal do $J$ ($J^{2}\neq-1$) & duas asserções falham \\
a estaca sem ponto fixo & tudo cai do mesmo lado --- \textbf{não termina} \\
\bottomrule
\end{tabular}
\end{center}

\noindent \textbf{Duas das três não dão resultado errado: não dão resultado nenhum.} É o teste mais
forte de que a lei não é decoração --- \emph{sem ela não há algoritmo}.

\section{A cruz, e por que quatro bastam}

Ordenar \emph{pares} muda o que se pode dizer, e a razão é a torre: em $\mathbb{R}^{n}$ a ordem já
está na estrutura. Num número solto há uma coordenada e a órbita das involuções \textbf{degenera};
num par há duas, e é aí que a bidualidade se vê.

\smallskip
\noindent \textbf{A estaca move; a cruz mede o que não se moveu}, e tem duas coordenadas:
\[
  x\oplus x^{\dagger} \;=\; 2c \quad\text{(invariante --- \emph{a que mede})},
  \qquad
  x\otimes x^{\dagger} \quad\text{(varia, e é máxima no ponto fixo --- \emph{a que ordena})}.
\]

\noindent A soma não distingue nada: vale $2c$ em todo o conjunto. O produto satura exatamente em
$x=c$, e \textbf{o ponto fixo é essa saturação} --- a passagem da dimensão. \emph{A cruz lança a
projeção no dual}, e é ela que dá a segunda involução.

\smallskip
\noindent \textbf{Com pares, as duas são independentes e a órbita dá quatro saídas distintas.}
Quatro bastam; não foi preciso derivar oito.

\medido{a soma $x\oplus x^{\dagger}$ é invariante em $11$ pontos e vale $2c$; o produto é máximo
em $x=c$; e ordenar pares com uma estaca por coordenada dá $4$ saídas distintas de $4$ --- com a
mutação (as duas estacas iguais) a cair para $2$, que é a degenerescência (\code{tests/dualsort.c}).}

\smallskip
\noindent \emph{Uma tentativa falhada, registada porque poupa a próxima.} Antes de olhar para a
cruz, procurei a segunda involução \emph{inventando-a}: propus três (ordem, sentido, lado) e medi
--- das \textbf{oito} combinações saíam \textbf{duas} saídas, porque as três dependiam do produto
dos sinais e essa dependência fora eu que a construíra no código. \textbf{A segunda involução não
se inventa: já está escrita, e é a cruz.}

\section{Comparado com os outros}

Até aqui nada se comparou: o algoritmo foi construído das suas peças e medido contra si próprio.
Comparar vem agora, e sem escolher o terreno --- as distribuições incluem \emph{de propósito} o
caso mau para quem parte por valor, que é o caso mau deste algoritmo.

\begin{center}\small
\begin{tabular}{@{}lrrrrr@{}}
\toprule
$n=10^{6}$, tempo em ms & \textbf{Dual Sort} & \code{qsort} & quicksort & mergesort & heapsort \\
\midrule
uniforme            & $89{,}2$ & $99{,}1$ & $\mathbf{58{,}2}$ & $76{,}6$ & $107{,}9$ \\
quase ordenado      & $46{,}2$ & $34{,}6$ & $\mathbf{19{,}0}$ & $27{,}7$ & $54{,}8$ \\
invertido           & $25{,}3$ & $28{,}6$ & $\mathbf{7{,}3}$  & $17{,}6$ & $54{,}5$ \\
\textbf{8 valores distintos} & $\mathbf{10{,}7}$ & $50{,}8$ & $20{,}0$ & $34{,}6$ & $49{,}4$ \\
enviesado           & $43{,}3$ & $66{,}8$ & $\mathbf{27{,}8}$ & $47{,}3$ & $72{,}6$ \\
\bottomrule
\end{tabular}
\end{center}

\noindent \textbf{Onde ele ganha, ganha por muito, e a razão é a fronteira.} Com poucos valores
distintos faz $10{,}7$\,ms contra $20{,}0$ do quicksort e $50{,}8$ do \code{qsort} da biblioteca ---
\emph{quase cinco vezes mais rápido} que este. O motivo é o \textbf{ponto fixo}: os elementos iguais
ao centro caem na fronteira, \textbf{não pertencem a nenhum lado, e não voltam a descer}. Uma
partição de duas vias não tem onde os pôr e re-particiona-os em cada nível.

\smallskip
\noindent \textbf{E onde perde, perde de forma consistente}, o que também tem de ser dito: o
quicksort com mediana de três é mais rápido em quatro dos cinco casos. Contra o \code{qsort} da
biblioteca, o Dual Sort ganha em quatro dos cinco.

\smallskip
\noindent \emph{A leitura honesta é esta}: não é o mais rápido em geral, é o mais rápido
\textbf{onde há repetição} --- e isso não é um acaso de implementação, é o trial a fazer o seu
trabalho. Três medições independentes dão dispersão abaixo de $2\%$.

\medido{$n=10^{6}$, três corridas, dispersão ${<}2\%$; a saída é conferida crescente em cada caso
(\code{tools/bench\_dualsort.c}).}

\section{Uma nota sobre o percurso}

A primeira versão deste trabalho não usava peça nenhuma desta teoria --- sem estaca, sem Lei 2, sem
bidualidade --- e terminava a justificar-se com algoritmos de fora. \textbf{Era o resultado a pedir
licença a uma teoria que não é a sua}, e foi reescrito.

\smallskip
\noindent O segundo erro foi da mesma família, e apareceu como falha de execução: a partição
mandava o \emph{ponto fixo} para um dos lados, deixava de reduzir, e a recursão não terminava.
\textbf{O defeito foi programar contra a lei que o próprio texto enuncia} --- que o ponto fixo não
pertence a nenhum lado. São três, e não dois.

\vfill
{\footnotesize\color{cinza}\noindent
Teoria, catálogo e medidores em
\href{https://goldenkingdom.patriatechnology.com}{goldenkingdom.patriatechnology.com}.
CC BY 4.0 (textos) $\cdot$ MIT (código).\par}

\end{document}
